\documentclass[12pt,a4paper]{article}

\usepackage{amsmath, amssymb, amsthm}
\usepackage{geometry}
\geometry{margin=2.5cm}

\title{Modelarea Matematică a Sistemelor Materiale \\ \large Principii Fundamentale Folosite în Modelare}
\author{}
\date{}

\begin{document}

\maketitle

\section{Principii fundamentale folosite în modelare}

Modelarea matematică a sistemelor materiale se bazează pe un set de principii fundamentale care guvernează comportamentul fizic al materialelor și proceselor. Aceste principii sunt esențiale pentru formularea ecuațiilor care descriu evoluția sistemului.

\subsection{4.1. Legi de conservare}

Legile de conservare exprimă faptul că anumite mărimi fizice nu pot fi create sau distruse, ci doar transformate sau transferate. Cele mai importante sunt:

\begin{itemize}
    \item \textbf{Conservarea masei}
    \item \textbf{Conservarea energiei}
    \item \textbf{Conservarea impulsului}
    \item \textbf{Conservarea sarcinii electrice}
\end{itemize}

Aceste principii conduc la ecuații generale de bilanț de forma:


\[
\frac{d}{dt}(\text{cantitate}) = \text{intrări} - \text{ieșiri} + \text{producție}.
\]



În formă diferențială locală, o lege de conservare poate fi exprimată ca:


\[
\frac{\partial u}{\partial t} + \nabla \cdot \mathbf{J} = S,
\]


unde:
\begin{itemize}
    \item $u$ este densitatea mărimii conservate,
    \item $\mathbf{J}$ este fluxul acesteia,
    \item $S$ este termenul de producție sau sursă.
\end{itemize}

Această formulare stă la baza ecuațiilor de transport, difuzie, conducție termică, mecanică fluidelor și electromagnetism.

\subsection{4.2. Relații constitutive}

Relațiile constitutive descriu modul în care materialul răspunde la stimuli externi. Ele leagă variabilele interne ale materialului (tensiuni, deformații, fluxuri, câmpuri) de variabilele externe (forțe, temperatură, viteze, potențiale).

Exemple fundamentale:

\begin{itemize}
    \item \textbf{Legea lui Hooke (elasticitate liniară)}  
    

\[
    \sigma = E \, \varepsilon,
    \]


    unde $\sigma$ este tensiunea, $\varepsilon$ este deformația, iar $E$ este modulul lui Young.

    \item \textbf{Legea lui Fourier (conducție termică)}  
    

\[
    \mathbf{q} = -k \nabla T,
    \]


    unde $\mathbf{q}$ este fluxul de căldură, $k$ conductivitatea termică, iar $T$ temperatura.

    \item \textbf{Legea lui Newton (vâscozitate)}  
    

\[
    \tau = \mu \frac{du}{dy},
    \]


    unde $\tau$ este tensiunea tangențială, $\mu$ vâscozitatea, iar $u$ viteza fluidului.

    \item \textbf{Legea lui Ohm (conducție electrică)}  
    

\[
    \mathbf{J} = \sigma \mathbf{E},
    \]


    unde $\mathbf{J}$ este densitatea de curent, $\sigma$ conductivitatea electrică, iar $\mathbf{E}$ câmpul electric.
\end{itemize}

Relațiile constitutive sunt esențiale pentru închiderea sistemelor de ecuații provenite din legile de conservare.

\subsection{4.3. Ecuații diferențiale}

Majoritatea sistemelor materiale sunt descrise prin ecuații diferențiale, deoarece acestea surprind evoluția în timp și spațiu a variabilelor fizice.

\subsubsection*{Ecuații diferențiale ordinare (ODE)}
Apar în sisteme cu un număr finit de variabile de stare:


\[
\frac{dx}{dt} = f(x,t).
\]



\subsubsection*{Ecuații diferențiale parțiale (PDE)}
Apar în fenomene distribuite spațial:


\[
\frac{\partial u}{\partial t} = \mathcal{L}(u, \nabla u, \nabla^2 u, \ldots),
\]


unde $\mathcal{L}$ este un operator diferențial.

Exemple:
\begin{itemize}
    \item Ecuația căldurii:
    

\[
    \frac{\partial T}{\partial t} = \alpha \nabla^2 T.
    \]


    \item Ecuația undelor:
    

\[
    \frac{\partial^2 u}{\partial t^2} = c^2 \nabla^2 u.
    \]


    \item Ecuațiile Navier–Stokes:
    

\[
    \rho \left( \frac{\partial \mathbf{v}}{\partial t} + \mathbf{v}\cdot\nabla \mathbf{v} \right)
    = -\nabla p + \mu \nabla^2 \mathbf{v}.
    \]


\end{itemize}

\subsubsection*{Sisteme diferențiale neliniare}

Modelele neliniare apar în:
\begin{itemize}
    \item dinamica fluidelor;
    \item elasticitate neliniară;
    \item reacții chimice;
    \item fenomene cu feedback puternic.
\end{itemize}

Acestea pot prezenta:
\begin{itemize}
    \item bifurcații,
    \item instabilități,
    \item comportament haotic.
\end{itemize}

\end{document}